\chapter{第九届大唐杯A组省赛}
\fancyhead[L]{\color{H6}\kaishu\faIcon{atom}\;第九届大唐杯A组省赛}
\fancyhead[R]{\color{H6}\kaishu\rightmark\,}

\date{2022年4月10日}{}{\href{https://github.com/xubenshan/dtcup}{\textbf{才疏学浅的小熊}}}
{}
{https://www.xiaohongshu.com/user/profile/6192219d0000000021028647}{小红书}
{https://space.bilibili.com/692546609}{B站账号}




\section{单选题（共30小题，共90分）}



\begin{choice}{C}[]
    5G NR 结构中，1 个子帧的时隙数不可能有
    \begin{tasks}(4)
        \task 2
        \task  1
        \task 3
        \task 4
    \end{tasks}
\end{choice}


\begin{choice}{D}[]
    NR 系统中，下列哪项不属于 5G 无线技术
    \begin{tasks}(4)
        \task 大规模天线阵列
        \task 超密集组网
        \task 新型多址
        \task 网络功能虚拟化
    \end{tasks}
\end{choice}




\begin{choice}{D}[]
    在 5G 网络架构中，我们所提到的软件和硬件解耦主要通过下面哪种技术实现
    \begin{tasks}(4)
        \task 移动边缘计算
        \task 切片技术
        \task SDN
        \task NFV
    \end{tasks}
\end{choice}




\begin{choice}{D}[]
    5G NR 系统 R15 版本中PDSCH 支持最大调制方式为
    \begin{tasks}(4)
        \task 512QAM
        \task 128QAM
        \task 64QAM
        \task 256QAM
    \end{tasks}
\end{choice}


\begin{choice}{D}[]
    对于大唐 5G 设备，单模集束光纤在 BBU 与 AAU 路由长度超过多少米使用
    \begin{tasks}(4)
        \task 40
        \task 60
        \task 80
        \task 100
    \end{tasks}
\end{choice}

\begin{choice}{D}[]
    5G NR 系统中基于基站间 Xn 链路切换，源基站侧向目标基站发送第一条信令为
    \begin{tasks}(2)
        \task SN Status Transfer
        \task Handover Required
        \task Handover Request Acknowledge
        \task Handover Request
    \end{tasks}
\end{choice}




\begin{choice}{A}[]
    5G系统帧结构从时间的维度上对资源进行划分和使用，无线接口共多少个无线帧循环
    \begin{tasks}(4)

        \task 1024
        \task 4096
        \task 612
        \task 2048
    \end{tasks}
\end{choice}


\begin{choice}{D}[]
    以下不属于线性回归模型效果判断评估项的是
    \begin{tasks}(4)
        \task R2
        \task RMSE
        \task MSE
        \task F1 值
    \end{tasks}
\end{choice}


\begin{choice}{B}[]
    5G NR 系统中，5G 基站和核心网通过哪个接口相连
    \begin{tasks}(4)
        \task Xn
        \task NG
        \task X2
        \task S1
    \end{tasks}
\end{choice}


\begin{choice}{A}[]
    5G NR 系统 RRC 重建流程，若 UE 上下文不能被恢复但有资源，基站会触发哪条信令
    \begin{tasks}(2)
        \task RRC Setup
        \task 不回任何消息
        \task RRCestablishment
        \task RRC Reject
    \end{tasks}
\end{choice}

\begin{choice}{A}[]
    目前，我国 5G 商用网采用的 NSA 组网部署方式为哪一种
    \begin{tasks}(4)
        \task option3x
        \task option7x
        \task option7
        \task option3
    \end{tasks}
\end{choice}


\begin{choice}{A}[]
    以下选项中哪个不是5G下行物理信道
    \begin{tasks}(2)
        \task PCFICH
        \task PBCH
        \task PDCCH
        \task PDSCH
    \end{tasks}
\end{choice}


\begin{choice}{B}[]
    在 5G NSA 模式 opion3x EN-DC 中，下面哪个信令可以携带 SCG Add 或 PSCell Change 相关 SCG 配置
    \begin{tasks}(1)
        \task NR 侧的 RRCReconfiguration
        \task LTE 侧的 RRCConnectionReconfiguration
        \task  LTE 侧的 RRCConnectionSetup
        \task NR 侧的 RRCConnectionSetup
    \end{tasks}
\end{choice}


\begin{choice}{C}[]
    5G NR 系统中，PDCCH信道采用 Polar 码信道编码方式，调制方式为
    \begin{tasks}(2)
        \task 16QAM
        \task 256QAM
        \task QPSK
        \task 64QAM
    \end{tasks}
\end{choice}

\begin{choice}{D}[]
    NR 系统中，全双工的最大难点在于
    \begin{tasks}(4)
        \task 调制
        \task 解调
        \task 滤波
        \task 干扰抑制
    \end{tasks}
\end{choice}

\begin{choice}{B}[]
    车联网是实现自动驾驶的必要条件，以下不属于车联网优点的为
    \begin{tasks}(2)
        \task 可提供丰富的网联应用
        \task 可增加感知范围，但感知成本会相应提升
        \task 提升交通效率
        \task 提升安全性
    \end{tasks}
\end{choice}



\begin{choice}{D}[]
    5G 系统帧结构，哪种 SCS 在常规 CP 和扩展 CP 时均可采用
    \begin{tasks}(4)
        \task 120
        \task 30
        \task 15
        \task 60
    \end{tasks}
\end{choice}

\begin{choice}{A}[]
    在 5G 网络架构中，以下选项哪一项是 AMF 的功能
    \begin{tasks}(4)
        \task 注册管理
        \task 下行数据的通知
        \task 无线资源分配
        \task 会话的建立修改删除
    \end{tasks}
\end{choice}


\begin{choice}{A}[]
    在 5G NR 超密集组网的场景中，哪一项是解决边界效应的关键
    \begin{tasks}(4)
        \task 小区虚拟化
        \task 软扇区
        \task 回传技术
        \task 全频谱接入
    \end{tasks}
\end{choice}



\begin{choice}{A}[]
    5G NR 通信系统 R15 中，CORESET在时域上最多可以使用几个 OFDM 符号
    \begin{tasks}(4)
        \task 3
        \task 2
        \task 1
        \task 4
    \end{tasks}
\end{choice}



\begin{choice}{B}[]
    下列哪个因素不影响 AAU 正常接入
    \begin{tasks}(2)
        \task DCPD 给 AAU 供电的开关未闭合
        \task 修改了小区物理 ID 列表
        \task 对应的光模块插反了
        \task 对应的 HBPOD 板卡状态为未初始化状态
    \end{tasks}
\end{choice}




\begin{choice}{A}[]
    5G NR，UE 在空闲状态检测到的 RSRP 信号强度是通过测量下面哪个信号得到
    \begin{tasks}(4)
        \task PBCH DMRS
        \task  SSS
        \task PSS
        \task CSI-RS
    \end{tasks}
\end{choice}

\begin{choice}{B}[]
    大唐 5G 基站设备，BBU 接地线采用多少平方黄绿接地线
    \begin{tasks}(4)
        \task 6
        \task 16
        \task 35
        \task 25
    \end{tasks}
\end{choice}

\begin{choice}{B}[]
    在 5G NR 网络中终端在初始接入过程中，基站向核心网发的第一条信息是
    \begin{tasks}(2)
        \task UE Capability Information
        \task Initial UE Message
        \task  RRC Setup Complete
        \task  Uplink NAS Transfer
    \end{tasks}
\end{choice}


\begin{choice}{C}[]
    5G NR 系统中，RRC INACTVE 状态，如果要恢复业务，则会触发哪条信令
    \begin{tasks}(1)
        \task RRC connection reject
        \task RRC connection reconfiguration
        \task RRC resume request
        \task  RRC connection request
    \end{tasks}
\end{choice}


\begin{choice}{A}[]
    5G NR 中哪类 SRB 仅承载 NAS 消息，映射到 DCCH 信道
    \begin{tasks}(4)
        \task SRB2
        \task SRB0
        \task SRB3
        \task SRB1
    \end{tasks}
\end{choice}



\begin{choice}{B}[]
    5G NR 系统中，MIB 一个最重要作用是通知 UE 如何获取哪个消息
    \begin{tasks}(4)
        \task SIB3
        \task SIB1
        \task SIB2
        \task SIB4
    \end{tasks}
\end{choice}


\begin{choice}{A}[]
    在 5G NR 中 PBCH 信道上相邻DM-RS 在频域上的间隔是几个子载波
    \begin{tasks}(4)
        \task 4
        \task 2
        \task 6
        \task 8
    \end{tasks}
\end{choice}


\begin{choice}{C}[]
    5G 控制信道资源调度的基本单位 CCE 与 REG 之间的正确关系为
    \begin{tasks}(2)
        \task 1CCE =8RE
        \task 1CCE =9RB
        \task 1CCE =6REG
        \task 1CCE = 7RBG
    \end{tasks}
\end{choice}


\begin{choice}{D}[]
    C-V2X 中，下面有关 MEC 部署位置不正确的是
    \begin{tasks}(2)
        \task 与基站共址
        \task 接入汇聚机房
        \task 骨干汇聚机房
        \task 与 RRU 共址
    \end{tasks}
\end{choice}




\section{多选题（共20小题，共80分）}

\begin{choice}{\;ABC\;}[]
    NR 中测量报告特点包含哪些
    \begin{tasks}(1)
        \task 包含在测量报告中的小区和波束测量量由网络配置
        \task 网络配置的属于黑名单的小区不用于事件评估和上报，相反当网络配置白名单时，只有属于白名单的小区用于事件评估和上报
        \task 测量报告包括触发报告的相关测量配置的测量标识
        \task 要报告的非服务小区的数量不能通过网络配置来限制
    \end{tasks}
\end{choice}

\begin{choice}{\;ABCD\;}[]
    在 5G NR 中，RRC 连接建立完成后，对 UE 进行测量配置，根据 38.300 协议，测量类型包括
    \begin{tasks}(2)
        \task 基于 CSI-RS 的异频测量
        \task 基于 SSB 的同频测量
        \task 基于 CSI-RS 的同频测量
        \task 基于 SSB 的异频测量
    \end{tasks}
\end{choice}

\begin{choice}{\;ABC\;}[]
    关于越区覆盖，以下描述正确的是哪些项
    \begin{tasks}(1)
        \task 越区覆盖可能导致 UE 发射功率受限，引起接入失败
        \task 越区覆盖可能导致越区覆盖基站与周边基站负荷不均。周边基站资源利用率低
        \task 越区覆盖可能形成“孤岛”，导致掉话
        \task 对于 NR 系统而言越区覆盖不会对其邻近小区形成干扰
    \end{tasks}
\end{choice}

\begin{choice}{\;ABC\;}[]
    下面哪些接口属于 C−V2X 的 PC5 接口
    \begin{tasks}(2)
        \task 车之间
        \task 车与 RSU(UE 级别)
        \task 车与人
        \task 车与基站
    \end{tasks}
\end{choice}


\begin{choice}{\;ACD\;}[]
    5G NR 中，RRC 状态包括哪些
    \begin{tasks}(2)
        \task RRC CONNECTED
        \task  RM DEREGISTERED
        \task RRC IDLE
        \task RRC INACTIVE
    \end{tasks}
\end{choice}

\begin{choice}{\;ABCD\;}[]
    5G 的网络架构描述正确的是
    \begin{tasks}(1)
        \task 控制平面引入新的基于服务的设计思路
        \task 核心网抽取会话管理功能，形成独立模块 SMF
        \task 分布式的用户面功能 UPF，无需汇聚点，支持多会话
        \task 架构设计以网元功能为单位，支持网络切片
    \end{tasks}
\end{choice}

\begin{choice}{\;BD\;}[]
    5G NR 系统中针对系统消息部分采用了不同的方式，其中最小系统消息包括
    \begin{tasks}(4)
        \task SIB2
        \task MIB
        \task SIB3
        \task SIB1
    \end{tasks}
\end{choice}


\begin{choice}{\;ABCD\;}[]
    在 5G 网络优化数据分析中，属于机器学习流程的选项是
    \begin{tasks}(2)
        \task 数据收集
        \task 数据清洗
        \task 特征工程
        \task 数据建模
    \end{tasks}
\end{choice}


\begin{choice}{\;ABC\;}[]
    在5G所支持的SCS中，能被所有物理信道和信号使用的SCS是
    \begin{tasks}(4)
        \task 120KHz
        \task 30KHz
        \task 15KHz
        \task 60KHz
    \end{tasks}
\end{choice}


\begin{choice}{\;BC\;}[]
    C-V2X 可提供两种通信接口，分别是
    \begin{tasks}(4)
        \task Xn
        \task PC5
        \task Uu
        \task Ng
    \end{tasks}
\end{choice}


\begin{choice}{\;AD\;}[]
    下面有关网络切片说法正确的是
    \begin{tasks}(1)
        \task 每个切片之间是逻辑上隔离
        \task 在一个硬件基础设施切分出多个物理的端到端网络
        \task 每个切片之间是物理上隔离
        \task 在一个硬件基础设施切分出多个虚拟的端到端网络
    \end{tasks}
\end{choice}

\begin{choice}{\;BCD\;}[]
    5G NR 系统中，SIB2 主要包含哪些内容
    \begin{tasks}(2)
        \task 同频小区专用重选消息
        \task 小区重选系统内异频相关信息
        \task 小区异系统重选公共消息
        \task 小区重选系统内同频公共信息
    \end{tasks}
\end{choice}

\begin{choice}{\;ABD\;}[]
    针对大唐 5G 产品，进行巡检工作时，常用的工具有哪些
    \begin{tasks}(2)
        \task uem5000
        \task LMT
        \task  OSP
        \task NodeBSpider
    \end{tasks}
\end{choice}

\begin{choice}{\;BD\;}[]
    以下属于 AI 技术赋能 5G 方面应用是
    \begin{tasks}(4)
        \task 自然语言处理
        \task 资源调度
        \task 语音识别
        \task 网络优化
    \end{tasks}
\end{choice}

\begin{choice}{\;BD\;}[]
    下列数据类型，哪些属于 5G 物联网
    \begin{tasks}(2)
        \task 交互类数据
        \task 控制类数据
        \task  传输类数据
        \task 采集类数据
    \end{tasks}
\end{choice}

\begin{choice}{\;ABD\;}[]
    5GR15 版本，PDCCH 信道 CCE 聚合等级可能为
    \begin{tasks}(4)
        \task 4
        \task 2
        \task 3
        \task 1
    \end{tasks}
\end{choice}

\begin{choice}{\;AD\;}[]
    5G NR 系统中，UE 通过小区搜索过程完成下行同步，在此过程需要哪些信号参与
    \begin{tasks}(4)
        \task PSS
        \task CRS
        \task SRS
        \task SSS
    \end{tasks}
\end{choice}

\begin{choice}{\;ACD\;}[]
    5G NR 系统中，哪种情形下只能进行基于竞争的随机接入
    \begin{tasks}(1)
        \task 无线链路失败后进行初始接入
        \task 切换时进行随机接入
        \task 由 Idle 状态进行初始接入
        \task 在 Active 情况下，上行数据到达，如果没有建立上行同步，或者没有资源发送调度请求，则需要随机接入
    \end{tasks}
\end{choice}

\begin{choice}{\;ABD\;}[]
    在 5G 帧结构中单个时隙内的 OFDM 符号类型有哪些
    \begin{tasks}(4)
        \task downlink
        \task flexible
        \task empty
        \task uplink
    \end{tasks}
\end{choice}

\begin{choice}{\;ABD\;}[]
    5G 网络部署 Massive MIMO，主要用于解决哪些问题
    \begin{tasks}(2)
        \task 大容量问题
        \task 深度覆盖问题
        \task 穿透损耗问题
        \task 天面不足问题
    \end{tasks}
\end{choice}



\section{判断题（共10小题，共30分）}

\begin{choice}{\;正确\;}[]
    5G NR 系统 RRC 重建流程，若没有无线资源，此时基站会触发拒绝消息。
\end{choice}


\begin{choice}{\;错误\;}[]
    5G 系统中，只有信道有调制方式，信号是没有调制方式。
\end{choice}

\begin{choice}{\;正确\;}[]
    5G 是 AI 落地的新动能，AI 将助推 5G 的更好发展。
\end{choice}

\begin{choice}{\;错误\;}[]
    国内运营商 5G 网络时钟源优先级最高的是北斗。
\end{choice}

\begin{choice}{\;正确\;}[]
    5G R15 版本小区带宽 sub-6GHz 包括:5/10/15/20/25/30/40/50/60/80/100MHz。
\end{choice}

\begin{choice}{\;错误\;}[]
    5G NR 帧结构，1 个时隙的符号数与 $\mu $ 值有关。

\end{choice}


\begin{choice}{\;错误\;}[]
    5G 解决决策问题，AI 解决通信问题。

\end{choice}

\begin{choice}{\;正确\;}[]
    5G 网络架构新增 NRF，可以实现网络功能服务管理的自动化(自动注册/更新/去注册动发现和选择、状态检测、认证授权等)。

\end{choice}

\begin{choice}{\;错误\;}[]

    5G NR 系统中，UE 在 RRC IDLE 态和 RRC INACTIVE 态的寻呼标识为 P-RNTI。
\end{choice}


\begin{choice}{\;错误\;}[]
    5G 系统中，每一个物理信道的数字调制方式是唯一的。

\end{choice}

